\documentclass[../../../include/open-logic-section]{subfiles}

\begin{document}

\olfileid[pt]{sth}{replacement}{ref}
\olsection{Substituição e Reflexão}

A nossa última tentativa de justificar a Substituição, via \stagesinex, começa com um resultado profundo e elegante:\footnote{Lembrete: todas as fórmulas podem ter parâmetros (salvo indicação explícita em contrário).}
%In this, I will use overlining, such as $x_1, \ldots, x_n$, to
%abbreviate ``$x_1, \ldots, x_n$'': 
\begin{thm}[Esquema de reflexão]\ollabel{reflectionschema}
Para qualquer fórmula $\phi$:
\[
\forall \alpha \exists \beta > \alpha (\forall x_1 \ldots, x_n \in
V_\beta)(\phi(x_1, \ldots, x_n) \liff \phi^{V_\beta}(x_1, \ldots, x_n))
\]
\end{thm}
\noindent 
Como na \olref[sth][replacement][strength]{formularelativization}, $\phi^{V_\beta}$ é o resultado de restringir todos os
quantificadores de $\phi$ ao conjunto~$V_\beta$. Assim, intuitivamente, a Reflexão
diz o seguinte: se $\phi$ é verdadeira em toda a hierarquia, então $\phi$ é
verdadeira em arbitrariamente muitos \emph{segmentos iniciais} da hierarquia. 

\citet{Montague1961} e \citet{Levy1960} mostraram que (formulações
adequadas da) Substituição e da Reflexão são equivalentes,
módulo~$\Z$, de modo que acrescentar qualquer uma delas dá~$\ZF$. (Provamos estes resultados na \olref[sth][replacement][refproofs]{sec}.) Dada esta
equivalência, poder-se-á esperar justificar a Reflexão e a Substituição via
\stagesinex{} do seguinte modo: dado \stagesinex, a hierarquia deverá ser
muito, muito alta; tão alta, de fato, que nada do que podemos dizer sobre ela é
suficiente para limitar a sua altura. E podemos entender isto como a
ideia de que, se qualquer sentença~$\phi$ é verdadeira em toda a hierarquia,
então é verdadeira em arbitrariamente muitos segmentos iniciais da hierarquia.
E isso é precisamente a Reflexão. 

De novo, isto parece uma tentativa genuinamente promissora de fornecer uma
justificação intrínseca da Substituição. Mas há demasiado a dizer sobre o assunto aqui. Cabe-lhe agora decidir por si se
tem êxito.\footnote{Poderá, no entanto, querer continuar lendo \citet[95--100]{Incurvati2020}.}

\end{document}
